\begin{abstract}
Off-policy evaluation for the current period must often rely on logs collected over many
earlier periods, with only a small sample from the period itself. The usual remedies for
non-stationarity, recency windows, decay and forecasting, treat every change in the logs
alike, so they discard data that is still informative. We argue that the useful question
is not how recent a period is but which mechanism changed: a change in the context
distribution or in the logging policy can be corrected by reweighting, while only a
change in the conditional reward biases a reweighted estimate. We propose O-DRIFT, which
segments the history with a per-action residual test that is invariant to context and
logger drift, fits joint reward models across segments, reuses older segments with a
drift correction, and pools the segment estimates at minimum variance with a
random-effects interval. In a pre-registered evaluation with a one-shot test split, O-DRIFT
outperforms a cross-domain estimator that groups periods by mean reward in all 15
reward-drift settings, and replacing its residual test by mean-reward grouping causes
spurious splits under context and logger drift. On semi-synthetic Open Bandit data,
however, drift-corrected reuse is strongly biased; we report this failure and show
why validation on the small current-period sample cannot repair it.
\keywords{Off-policy evaluation \and Non-stationarity \and Concept drift \and Recommender systems}
\end{abstract}
